% Generated draft: requires mathematical and source review.
\documentclass[11pt]{article}
\usepackage[margin=25mm]{geometry}
\usepackage{fontspec}
\setmainfont{Latin Modern Roman}
\usepackage{amsmath,amssymb,mathtools,unicode-math}
\setmathfont{Latin Modern Math}
\usepackage{enumitem,needspace,xurl,hyperref}
\setlength{\parindent}{0pt}
\setlength{\parskip}{5pt}
\setlength{\emergencystretch}{3em}
\newcommand{\N}{\mathbb{N}}
\newcommand{\Z}{\mathbb{Z}}
\newcommand{\Q}{\mathbb{Q}}
\newcommand{\R}{\mathbb{R}}
\newcommand{\C}{\mathbb{C}}
\newcommand{\F}{\mathbb{F}}
\newcommand{\A}{\mathbb{A}}
\newcommand{\PP}{\mathbb{P}}
\newcommand{\E}{\mathbb{E}}
\newcommand{\eps}{\varepsilon}
\newcommand{\dd}{\,\mathrm{d}}
\begin{document}
\section*{UnsolvedMath ID: 30003138}
\textbf{Source problem code:} catalogue entry
\subsection{Definitions and mathematical statement}
Use the conventional interpretation that a graph is finite, simple, and undirected; these conventions are implicit rather than explicit in the supplied record. An edge of a graph is \emph{nontriangular} if it belongs to no triangle. Write
\[
q(H)=\bigl|\{uv\in E(H):\text{there is no }w\in V(H)\setminus\{u,v\}\text{ with }uw,vw\in E(H)\}\bigr|.
\]
For nonnegative integers \(a,b,c\), let \(G(a,b,c)\) have a vertex partition \(A\sqcup B\sqcup C\), where \(|A|=a\), \(|B|=b\), and \(|C|=c\). The set \(A\) induces a clique, \(B\) and \(C\) are independent, every vertex of \(B\) is adjacent to every vertex of \(A\cup C\), and there are no edges between \(A\) and \(C\). Thus
\[
|E(G(a,b,c))|=\binom{a}{2}+ab+bc.
\]
\textbf{Conjecture (feasible-parameter formulation).} For every pair of integers \(n,e\) with \(n\geq 0\) and \(0\leq e\leq\binom{n}{2}\), there exist nonnegative integers \(a,b,c\) such that
\[
a+b+c=n,\qquad \binom{a}{2}+ab+bc\geq e,
\]
and, for every graph \(H\) with \(|V(H)|=n\) and \(|E(H)|\geq e\),
\[
q(H)\leq q(G(a,b,c)).
\]
The source says ``for every \(n\) and \(e\)'' without specifying their ranges. The displayed ranges make the optimization feasible under the graph conventions above; including \(n=0\) uses the empty-graph convention. The threshold is \emph{at least} \(e\) edges, not exactly \(e\) edges.

The supplied curation labels the problem partially solved and reports a sufficiently-large-\(n\) result of Gruslys--Letzter while retaining the unrestricted all-\(n\) formulation as open. This is a supplied literature assessment, not an independently verified theorem or status determination.
\subsection{Short English statement}
For every feasible vertex count and edge threshold, the conjecture asks whether some graph of the specified three-part form maximizes the number of edges lying in no triangle among all graphs with that vertex count and at least that many edges. The supplied dated curation reports a result for sufficiently large vertex counts and leaves the all-size formulation open.
\subsection{Sources}
\begin{itemize}
\item Supplied local UnsolvedMath registry record, permanent ID \(30003138\), title ``Extremal Graphs Maximizing Non-Triangular Edges,'' source problem code \texttt{catalogue entry}, in the Oberwolfach Reports Open Problems collection. The proposed year is recorded as 2016; no individual proposer is supplied.
\item Supplied report citation: \emph{Combinatorics and Probability (2017)}, Oberwolfach Reports, pp.~1189--1257, DOI \texttt{10.4171/owr/2016/22}. \url{https://doi.org/10.4171/owr/2016/22}. The record describes these reports as documenting workshop discussions and open problems.
\item Gruslys--Letzter, \emph{Minimising triangular edges}. \url{https://arxiv.org/abs/1605.00528}.
\item Letzter, described in the supplied record as a survey of extremal triangular-edge problems. \url{https://shohamlet.github.io/papers/triangular-edges.pdf}.
\item Literature assessment dated \texttt{2026-08-21}: the supplied curation reports that Gruslys--Letzter prove the exact structural conjecture for sufficiently large \(n\), and that later discussion explicitly leaves the version without a large-\(n\) assumption open. This is supplied dated curation derived from the OWR collection, not independently verified here.
\item The registry labels the statement \texttt{corrected\_verified} and states that it was checked against the official Oberwolfach report. This verification claim is supplied metadata, not a verification performed for this draft.
\end{itemize}
\end{document}
